The benchmark protocol follows the principle that estimator comparisons are meaningful only when data, specifications, starting values, and post-estimation definitions are aligned.
For each case, the validation layer records the public data source, alternatives, variables, availability filters, scaling transformations, initial parameter values, covariance definition, and reference backend.
When a backend exposes separate timing, the protocol reports parameter-estimation time and covariance/Hessian time separately.

The benchmark suite currently uses three groups of public data.
The first group contains public Biogeme datasets, including airline itinerary choice, parking choice, telephone service choice, Swissmetro, Optima, and London Passenger Mode Choice.
The second group contains R community datasets used by \texttt{mlogit}, including Fishing and ModeCanada.
The third group contains processed large-data candidates that are maintained as choice-set artifacts rather than raw survey dumps.
Small public datasets are committed to the repository, while large processed artifacts are distributed separately.

The numerical metrics are chosen to match the Apollo benchmark plan used to guide the project, but the paper tables report a single consistency flag rather than every raw difference.
A full-estimation case is marked consistent when all directly comparable quantities satisfy the following prespecified tolerances: absolute log-likelihood difference at most $10^{-4}$, maximum absolute parameter difference at most $5\times10^{-4}$, maximum probability difference at most $2\times10^{-4}$, and maximum standard-error difference at most $3\times10^{-2}$ when the same covariance definition is available.
Willingness-to-pay and elasticity checks use the same relative numerical scale when the model specification supports them.
For simulation-based models, the suite distinguishes full simulated maximum-likelihood estimation from fixed-parameter replay.
The Swissmetro mixed-logit full-estimation row uses a relaxed consistency rule against Biogeme under shared draws and a shared nonzero variance initial value, while simulation-based rows that are not yet fully estimated across backends use fixed-parameter replay with shared draws and mark consistency for the likelihood/probability kernel rather than for the complete simulated estimator.

The runtime metrics distinguish estimator work from reporting work.
Total runtime includes backend overhead when the benchmark script can measure it, while estimation time isolates parameter optimization and covariance time isolates Hessian or covariance calculations.
This split is important because some reference packages perform model construction, reporting, file output, or post-processing outside the core optimizer.
The paper-facing benchmark tables therefore report total runtime, while the finer estimation/covariance split remains available in the generated validation artifacts.
